\section{A finite bad-cycle bound on the primitive transversal}
\label{sec:finite-locality}

This part uses only the finitely many typed palettes of
Theorem~\ref{thm:joint-transversal}, SAP, and primitivity of each
individual sort. Write $P_{st}$ for the full pair palette and $*$
for strong endpoint composition: intermediate points must be distinct
from the endpoints and from one another in a realised finite path.
SAP allows us to replace repeated auxiliary points by fresh copies
while fixing their complete specified neighbourhoods. Consequently
strong products are associative and profiles of positive typed paths
can be computed one edge at a time. Nonempty compatible profiles have
nonempty products; every realised atomic neighbour set is infinite.

\begin{lemma}[Absorption by a full typed palette]
\label{lem:locality-absorb}
For every triple of owner sorts $s,t,u$, and every nonempty
atomic subsets $A\subseteq P_{st}$, $B\subseteq P_{tu}$,
not necessarily original exact menus, one has
\[
 A*P_{tu}=P_{su},\qquad P_{st}*B=P_{su}.
\]
In particular $A*P_{tt}=P_{st}$ and
$P_{st}*a=P_{su}$ for any atom $a\in P_{tu}$.
\end{lemma}
\begin{proof}
For the first equality, fix any target atom $c\in P_{su}$
and some $a\in A$.  Strongly amalgamate a realised $a$-edge
of sorts $(s,t)$ and a realised $c$-edge of sorts $(s,u)$
over their common $s$-endpoint.  The other endpoints remain
distinct, even if $t=u$, and their actual atom belongs to
$P_{tu}$.  This witnesses $c\in A*P_{tu}$.
For the second equality, fix any $c\in P_{su}$ and $b\in B$.
Strongly amalgamate a realised $c$-edge of sorts $(s,u)$
and a realised $b$-edge of sorts $(t,u)$ over the common
$u$-endpoint; their new $(s,t)$-edge lies in $P_{st}$ and
witnesses $c\in P_{st}*B$.  The reverse containments in
both cases are immediate from typing.  The two displayed
special cases follow by taking $u=t$ in the first formula
and $B=\{a\}$ in the second.
\end{proof}

\begin{lemma}[Every nonempty primitive loop palette eventually fills]
\label{lem:locality-positive}
On a primitive named sort $t$, let $\varnothing\ne W\subseteq
P_{tt}$ be an arbitrary atomic subset, with no assumption that
$W$ is an original exact menu or a represented word profile.
Put $m=|P_{tt}|\ge1$. Then
\[
                 W^{*2m}=P_{tt}
       \quad\text{and}\quad
                 W^{*n}=P_{tt}\quad(n\ge2m).
\]
Thus the eventual-fullness exponent is UNIFORM in the
nonempty atomic subset $W$, and the statement does not
merely assert eventual stabilization of its powers.
\end{lemma}
\begin{proof}
All original same-sort atomic orbitals on $t$ are symmetric,
so $W=\bar W$ and every $W^{*n}$ is symmetric.  For each
$a\in W$ choose a realised legal triangle with two $a$-edges,
whose third edge has some actual colour $b\in a*a$.
Strongly amalgamate two copies of this triangle over their
complete $b$-edge.  The off-base vertices stay distinct,
giving four distinct vertices on an $a$-coloured four-cycle.
Three consecutive $a$-edges thus realise the remaining $a$-edge
between the endpoints, and $a\in W^{*3}$.  Hence
$W\subseteq W^{*3}$ for arbitrary nonempty $W$, not just
word profiles.

Therefore the even powers
$E_k=W^{*2k}$ form an increasing chain of nonempty
subsets of the $m$-atom palette $P_{tt}$.
If $E_k=E_{k+1}$ for some $k$, all subsequent even
powers agree with $J=E_k$ by associativity, hence
$J*J=J$.  The symmetric invariant relation
$\Delta_t\cup J$ is then an ACTUAL equivalence:
the equality atom is added only in this interpreted
binary relation, never as a factor of the strong
off-diagonal product.  Since $J$ is nonempty,
primitivity forces $J=P_{tt}$.

Now consider $E_1,\ldots,E_m$.  If no adjacent
equality occurs, these are $m$ STRICTLY increasing
nonempty subsets of an $m$-element palette, so
$|E_m|\ge m$ and $E_m=P_{tt}$.
If any adjacent equality occurs, its stabilized
$J$ was already the full palette, and hence
$E_m=P_{tt}$ as well.  Thus in ALL cases
$W^{*2m}=P_{tt}$, including $m=1$.
By the full-palette absorption lemma
\cref{lem:locality-absorb},
multiplying the full profile by ANY nonempty
$W$ preserves fullness, so every higher power is
also full.  This proves the claimed explicit
uniform exponent without any STRONG identity atom.
\end{proof}

For a fixed starting sort $s$ define the \emph{proper-profile
automaton} to have states $(t,D)$, where
$\varnothing\ne D\subsetneq P_{st}$ is the profile of some positive
typed path, and edges
$(t,D)\to(u,D*a)$ whenever $a\in P_{tu}$ and the new profile
is again nonempty and proper. Full profiles are absorbing by
Lemma~\ref{lem:locality-absorb} and are omitted.

\begin{theorem}[Bounded cycle-obstruction locality]
\label{thm:locality-bound}
The proper-profile automaton is acyclic. There is an integer $\ell$
such that every positive typed word longer than $\ell$ has its full
endpoint palette as profile. One may take
\[
 \ell\le\max_s\sum_t(2^{|P_{st}|}-2).
\]
Every noncompletable simple typed cycle has length at most
$b=\max\{3,\ell+1\}$. In particular, by semantic homomorphic CCO,
a finite typed partial diagram can be completed whenever all its
specified cycles of length at most $b$ admit a typed completion.
\end{theorem}
\begin{proof}
Suppose a proper state $(t,D)$ recurs along a nonempty path through
the automaton. Let $W\subseteq P_{tt}$ be the positive closed-word
profile of this path. Then $D*W=D$. By associativity,
$D*W^{*n}=D$ for all $n$. By
Lemma~\ref{lem:locality-positive} the profile $W^{*n}$ is
$P_{tt}$ for sufficiently large $n$. But
$D*P_{tt}=P_{st}$ by Lemma~\ref{lem:locality-absorb}, contrary to
the properness of $D$. Therefore no state can repeat; a positive
path visiting proper states has at most the displayed number of
states. Once a full state is reached, its profile remains full.
This proves the word bound.

For a typed simple cycle of $m$ vertices, remove its closing edge.
The remaining path has length $m-1$; the specified cycle has a
homomorphic completion exactly when the closing-edge colour belongs
to that path's strong endpoint profile. The strong product computes
this profile with all auxiliary vertices distinct by SAP. If
$m-1>\ell$, the profile is full, so the cycle is always completable.
Thus any bad cycle has length at most $\ell+1$, and taking a maximum
with $3$ handles the conventional lower length bound for a simple
cycle. By typed semantic homomorphic CCO of the transversal, checking
all the cycles up to $b$ suffices for a completion of the full partial
diagram. The bound is chosen separately for each fixed finite owner
palette and need not be uniform over all quotient witnesses.
\end{proof}
